\documentclass[tikz,border=4pt]{standalone}
\usepackage{ctex}
\usepackage{amsmath}
\usepackage{pgfplots}
\pgfplotsset{compat=1.18}
\usetikzlibrary{calc, arrows.meta}

\begin{document}
% part14/02 抛物线+直线+动点综合图
% 抛物线 y = -x^2+2x+3; 交x轴于A(-1,0)B(3,0)；交y轴于C(0,3)
% 直线BC: y=-x+3; 动点P在弧BC上; 铅垂高示意
\begin{tikzpicture}
\begin{axis}[
  axis lines=center,
  xlabel={$x$}, ylabel={$y$},
  every axis x label/.style={at={(axis cs:4.2,0)}, anchor=west},
  every axis y label/.style={at={(axis cs:0,5.2)}, anchor=south},
  xmin=-2, xmax=4.5, ymin=-1, ymax=5.5,
  xtick={-2,-1,0,1,2,3,4},
  ytick={-1,0,1,2,3,4,5},
  tick label style={font=\small},
  width=7.5cm, height=7.5cm,
  thick,
]
  % 抛物线 y = -x^2+2x+3
  \addplot[blue, thick, domain=-1.2:3.2, samples=100] {-x^2+2*x+3};
  \node[blue, left, font=\small] at (axis cs:-0.5,3.5) {$y=-x^2+2x+3$};

  % 直线BC: y = -x+3
  \addplot[gray, dashed, domain=-0.3:3.3, samples=10] {-x+3};
  \node[gray, right, font=\small] at (axis cs:0.2,3.2) {$BC:\,y=-x+3$};

  % 关键点
  \fill (axis cs:-1,0) circle(2.5pt);
  \node[below left, font=\small] at (axis cs:-1,0) {$A(-1,0)$};

  \fill (axis cs:3,0) circle(2.5pt);
  \node[below right, font=\small] at (axis cs:3,0) {$B(3,0)$};

  \fill (axis cs:0,3) circle(2.5pt);
  \node[left, font=\small] at (axis cs:0,3) {$C(0,3)$};

  % 动点P（t=1.5时）
  \fill[red] (axis cs:1.5,3.75) circle(2.5pt);
  \node[above right, red, font=\small] at (axis cs:1.5,3.75) {$P(t,\,y_P)$};

  % Q：P在直线BC上的投影（铅垂高）
  \fill[gray] (axis cs:1.5,1.5) circle(2pt);
  \node[right, gray, font=\small] at (axis cs:1.5,1.5) {$Q$};
  \draw[gray, dashed] (axis cs:1.5,1.5) -- (axis cs:1.5,3.75)
    node[midway, right, gray, font=\small] {$h$};

  % 原点
  \node[below left, font=\small] at (axis cs:0,0) {$O$};

  % 顶点
  \fill (axis cs:1,4) circle(2pt);
  \node[above right, font=\small] at (axis cs:1,4) {顶点$(1,4)$};
\end{axis}
\end{tikzpicture}
\end{document}
